% Discussion.
\section{Discussion}
\label{sec:discussion}

\textbf{Geometry is the transferable quantity.} The central result is that under
a load shift the higher-order geometry of the reconstructed phase-space attractor
transfers markedly better than the raw waveform: the structural models and
GA-DCHG sit \GapCZeroThreeDynhg{}/\GapCZeroThreeGad{} and
\GapCZeroOneTwoThreeDynhg{}/\GapCZeroOneTwoThreeGad{} points above the 1D-CNN on
cwru\_03 and cwru\_0123 respectively. Raw signal statistics such as amplitude
envelopes and energy bands are strongly tied to the operating point, whereas the local
curvature, expansion, and filling of the attractor are more intrinsic to the
fault's dynamical signature. GA-DCHG realizes this by making the geometry active,
and attains the best accuracy on both protocols.

\textbf{Anchor selection as method, not assumption.} Because the invariant
geometry is \emph{protocol-dependent}, hard-coding a descriptor inherits the
shift of a domain-specific one whenever the regime changes. Screening by the
normalized $1$-Wasserstein distance is simple, uses only unlabeled target data,
and is a standalone contribution that should transfer to other phase-space-based
representations.

\textbf{Intervention audit and the ablation trade-off.} R4 perturbs the edge
geometry at test time, rather than only measuring attention--geometry
correlation: perturbing the edge geometry flips
\FlipScale{} of decisions on average, showing the decision depends on the geometry.
The component ablation (R3) is within seed noise on
cwru\_03 but is resolvable on cwru\_0123, where removing the anchor or the
temporal propagation degrades the model beyond its seed std. The geometry-space
alignment (C3) acts as a zero-accuracy-cost, explainable regularizer: on
cwru\_0123 removing it costs at most 0.010 accuracy points, on cwru\_03 it
sits within seed noise, and its physical interpretation is independent of any
accuracy gain.

\textbf{Limitations.} We evaluate only the CWRU load axis, which has four loads,
at three seeds, bounding the resolution of component-level differences; a more
fine-grained multi-rig or multi-seed study would sharpen the component estimates.
We also note that an MFPT cross-load transfer is not feasible in the same form:
its normal class is recorded at a single load, and an extreme low\,$\rightarrow$\,high
fault-load shift transferred for no model, including the 1D-CNN, so we do not
report an MFPT headline number. Finally, the anchor is selected from unlabeled
target data, legitimate in unsupervised domain adaptation but to be revisited if
no target data were available.

\textbf{Broader impact.} Because the anchor is a scalar geometry statistic, the
pipeline is lightweight and runs on CPU for these data sizes; the screening and
intervention audit give practitioners a concrete way to decide \emph{which}
geometry transfers for their rig and to verify a deployed model actually uses it.
The scalar anchor is also relevant to uncertainty-aware, active model
updating \cite{hong2025bayes,song2025quad}, which can attach probabilistic
guarantees as operating conditions evolve.
